\documentclass[a4paper,14pt]{extarticle}

% Базовые пакеты
\usepackage[utf8]{inputenc}
\usepackage[T2A]{fontenc}
\usepackage[russian]{babel}

% Настройка полей по ГОСТ (Левое 30мм, Правое 15мм, Верх и Низ 20мм)
\usepackage[left=3cm, right=1.5cm, top=2cm, bottom=2cm]{geometry}

% Пакеты для графики, таблиц и кода
\usepackage{graphicx}
\usepackage{tabularx}
\usepackage{listings}
\usepackage{hyperref}
\usepackage{setspace}

% Полуторный интервал
\onehalfspacing

\begin{document}

% =================================================================
% ТИТУЛЬНЫЙ ЛИСТ
% =================================================================
\begin{titlepage}
\begin{center}
\small
МИНИСТЕРСТВО СЕЛЬСКОГО ХОЗЯЙСТВА РОССИЙСКОЙ ФЕДЕРАЦИИ\\
ФЕДЕРАЛЬНОЕ ГОСУДАРСТВЕННОЕ БЮДЖЕТНОЕ ОБРАЗОВАТЕЛЬНОЕ УЧРЕЖДЕНИЕ ВЫСШЕГО ОБРАЗОВАНИЯ\\
«РОССИЙСКИЙ ГОСУДАРСТВЕННЫЙ АГРАРНЫЙ УНИВЕРСИТЕТ – МСХА имени К.А. ТИМИРЯЗЕВА»\\
Институт экономики и управления\\
Кафедра статистики и кибернетики\\[2cm]

\large
\textbf{КУРСОВОЙ ПРОЕКТ}\\
\normalsize
по дисциплине: Методы и средства проектирования информационных систем\\
на тему: \textbf{Проектирование и разработка веб-приложения «Персональный блог»}\\[2cm]
\end{center}

\begin{flushright}
Выполнил:\\
Студентка 3 курса группы ДЭ-317 Иванов И.И.\\[0.8cm]
Руководитель:\\
\underline{\hspace{5cm}}\\
(ученая степень, ученое звание, ФИО)\\[0.8cm]
Оценка \underline{\hspace{4cm}}\\
Дата защиты \underline{\hspace{3.4cm}}
\end{flushright}

\vfill
\begin{center}
Москва, 2023
\end{center}
\end{titlepage}

% =================================================================
% ЗАДАНИЕ НА КП
% =================================================================
\newpage
\begin{center}
\textbf{ЗАДАНИЕ\\НА КУРСОВОЙ ПРОЕКТ (КП)}
\end{center}
\vspace{0.5cm}
\noindent Обучающийся \underline{\hspace{12.5cm}}\\[0.3cm]
Тема КП \underline{\hspace{13.7cm}}\\[0.3cm]
Исходные данные к работе \underline{\hspace{10.2cm}}\\
\underline{\hspace{16.5cm}}\\[0.3cm]
Перечень подлежащих разработке в работе вопросов: \underline{\hspace{5.5cm}}\\
\underline{\hspace{16.5cm}}\\[0.3cm]
Перечень дополнительного материала \underline{\hspace{8.2cm}}\\
\underline{\hspace{16.5cm}}\\[1cm]
Дата выдачи задания «\underline{\hspace{1cm}}» \underline{\hspace{4cm}} 201\underline{\hspace{0.5cm}}г.\\[0.5cm]
Руководитель \underline{\hspace{7cm}}\\[0.5cm]
Задание принял к исполнению \underline{\hspace{5cm}}

% =================================================================
% РЕЦЕНЗИЯ
% =================================================================
\newpage
\begin{center}
\textbf{РЕЦЕНЗИЯ}\\
на курсовую работу/проект обучающегося
\end{center}
\vspace{0.5cm}
\noindent Обучающийся \underline{\hspace{12.5cm}}\\[0.3cm]
Учебная дисциплина \underline{\hspace{11.3cm}}\\[0.3cm]
Тема курсовой работы/проекта \underline{\hspace{9.5cm}}\\
\underline{\hspace{16.5cm}}\\[0.3cm]
Полнота раскрытия темы: \underline{\hspace{11cm}}\\
\underline{\hspace{16.5cm}}\\[0.3cm]
Оформление: \underline{\hspace{13.4cm}}\\
\underline{\hspace{16.5cm}}\\[0.3cm]
Замечания: \underline{\hspace{13.8cm}}\\
\underline{\hspace{16.5cm}}\\[0.5cm]
Курсовой проект отвечает предъявляемым к ней требованиям и заслуживает \underline{\hspace{4cm}} оценки.\\[0.8cm]
Рецензент \underline{\hspace{8cm}}\\[0.8cm]
Дата: «\underline{\hspace{1cm}}» \underline{\hspace{3cm}} 20\underline{\hspace{0.5cm}} г. \hfill Подпись: \underline{\hspace{3cm}}

% =================================================================
% АННОТАЦИЯ
% =================================================================
\newpage
\section*{Аннотация}
В ходе работы над курсовым проектом проведено теоретическое описание алгоритмов поиска, их назначение, виды, особенности реализации. По итогам проектирования реализовано пять алгоритмов поиска для такой структуры данных как массив. Инструментом реализации является высокоуровневый язык программирования Python (версия 3.0).

\textbf{Ключевые слова:} алгоритм, структуры данных, алгоритм поиска, программирование.

Проект состоит из 36 страниц и 5 приложений, использовано 18 литературных источников.

% =================================================================
% СОДЕРЖАНИЕ
% =================================================================
\newpage
\tableofcontents

% =================================================================
% ВВЕДЕНИЕ
% =================================================================
\newpage
\section*{Введение}
\addcontentsline{toc}{section}{Введение}
В современном информационном обществе персональный блог является не только средством самовыражения, но и важным инструментом для формирования профессионального портфолио, обмена опытом и накопления базы знаний. 

Использование современных технологий веб-разработки, таких как связка React и Express.js, позволяет создавать высокопроизводительные, масштабируемые и интуитивно понятные для пользователя веб-приложения. Разработка собственного решения, в отличие от использования готовых CMS-систем, предоставляет полный контроль над архитектурными слоями, механизмами безопасности, схемами данных и пользовательским интерфейсом. 

Написание такого проекта на базе языка строгого программирования TypeScript существенно повышает типобезопасность системы, снижает вероятность возникновения ошибок на этапе компиляции, что делает исследование методов построения подобных систем актуальной задачей для начинающих инженеров-программистов.

\textbf{Целью курсового проектирования} являются проектирование и программная реализация полноценного клиент-серверного веб-приложения «Персональный блог» с использованием современных JavaScript/TypeScript фреймворков, библиотек и архитектурных методологий.

Для достижения поставленной цели необходимо решить следующие задачи:
\begin{itemize}
    \item Провести детальный анализ предметной области, выявить основные функциональные требования к персональным блогам и рассмотреть существующие аналоги.
    \item Выбрать и научно обосновать технологический стек для реализации клиентской (Frontend) и серверной (Backend) частей приложения.
    \item Спроектировать реляционную структуру базы данных и описать схемы хранения пользовательских сущностей и контента.
    \item Разработать серверную часть (REST API) с использованием платформы Node.js и фреймворка Express.js.
    \item Разработать интерактивный пользовательский интерфейс на базе библиотеки React 19 с применением стейт-менеджера Zustand и методологии Feature-Sliced Design.
    \item Настроить процессы контейнеризации приложения с помощью Docker для упрощения последующего развертывания системы.
\end{itemize}

\textbf{Объектом исследования} являются современные технологии, архитектурные паттерны и методы разработки веб-приложений (клиент-серверные системы, Single Page Applications).

\textbf{Предметом исследования} являются процесс проектирования, разработки, типизации и развертывания веб-приложения «Персональный блог» на базе стека TypeScript, React, Express.js и Postgres.

% =================================================================
% ГЛАВА 1
% =================================================================
\newpage
\section{Теоретические основы разработки веб-приложений типа «Персональный блог»}

\subsection{Анализ предметной области и обзор существующих решений}
Здесь приводится понятие алгоритмов, их основное назначение, сферы применения. Также здесь можно привести виды алгоритмов поиска. Применение алгоритмов поиска в анализе и обработке данных. Текст. Текст. Текст [1].

Номер источника литературы, указанный в квадратных скобках, должен точно совпадать с номером источника в библиографическом списке.

\subsection{Теоретические основы построения современных веб-приложений (SPA, клиент-серверная архитектура)}
Описываем данный вид алгоритма. Пример оформления таблиц, рисунков, схем смотрите в методических указаниях. Нумерация таблиц, рисунков, листинга кода и так далее – сквозная.

\subsection{Особенности применения REST API при проектировании блога}
Описываем алгоритм линейного поиска. Суть алгоритма заключается в том, чтобы перебрать массив и вернуть индекс первого вхождения элемента, когда он найден.

% =================================================================
% ГЛАВА 2
% =================================================================
\newpage
\section{Практическая реализация веб-приложения «Персональный блог»}

\subsection{Обоснование выбора и описание инструментария реализации веб-приложения}
Данный параграф приводится в том случае, если для реализации алгоритмов выбран конкретный язык программирования или псевдокод. Например, если выбран язык программирования Python, то описывается назначение этого языка, его особенности, сравнение с другими языками программирования.

\subsection{Проектирование архитектуры приложения на основе методологии Feature-Sliced Design и структуры базы данных}
Здесь приводится конкретный пример работы линейного поиска и его поэтапная реализация простым описанием, псевдокодом, языком программирования Python или иным языком программирования.

\subsection{Разработка серверной части приложения на Express.js}
Текст. Текст. Текст. Текст. Текст.

\subsection{Разработка клиентской части приложения на React и Zustand}
Текст. Текст. Текст. Текст. Текст.

\subsection{Контейнеризация разработанного приложения с помощью Docker}
Этот раздел не является обязательным.

% =================================================================
% ЗАКЛЮЧЕНИЕ
% =================================================================
\newpage
\section*{Заключение}
\addcontentsline{toc}{section}{Заключение}
В заключении должно содержаться:
\begin{itemize}
    \item Краткое описание темы проекта (о чем проект).
    \item Вводный текст о достижении цели и решенных задачах курсового проекта.
    \item Выводы по каждой главе (что было сделано в каждой главе).
    \item Подтверждение актуальности и значимости поставленной задачи.
    \item Выявление направлений дальнейшего развития тематики курсового проекта.
\end{itemize}

% =================================================================
% БИБЛИОГРАФИЧЕСКИЙ СПИСОК
% =================================================================
\newpage
\section*{Библиографический список}
\addcontentsline{toc}{section}{Библиографический список}
\begin{enumerate}
    \item Кудрина, Е. В. Основы алгоритмизации и программирования на языке C\# : учебное пособие для вузов / Е. В. Кудрина, М. В. Огнева. — Москва : Издательство Юрайт, 2023. — 322 с.
    \item Трофимов, В. В. Алгоритмизация и программирование : учебник для вузов / В. В. Трофимов, Т. А. Павловская ; под редакцией В. В. Трофимова. — Москва : Издательство Юрайт, 2023. — 137 с.
    \item Чернышев, С. А. Основы программирования на Python : учебное пособие для вузов / С. А. Чернышев. — Москва : Издательство Юрайт, 2023. — 286 с.
\end{enumerate}

% =================================================================
% ПРИЛОЖЕНИЯ
% =================================================================
\newpage
\section*{Приложения}
\addcontentsline{toc}{section}{Приложения}

\textbf{Приложение 1}\\
Листинг 1.1 – Реализация алгоритма бинарного поиска в Python
\begin{lstlisting}[language=Python]
def binary_search(list1, low, high, n):
    if low <= high:
        mid = (high + low) // 2
        if list1[mid] == n:
            return mid
        elif list1[mid] > n:
            return binary_search(list1, low, mid - 1, n)
        else:
            return binary_search(list1, mid + 1, high, n)
    else:
        return -1
\end{lstlisting}

\vspace{1cm}
\textbf{Приложение 2. Инструментарий проекта (на основе README.md)}\\
На данный момент проект представляет собой CRUD-приложение для работы со списком задач (авторизация, создание, чтение, обновление и удаление задач, Drag-And-Drop).

\begin{table}[h]
\centering
\begin{tabularx}{\textwidth}{|l|X|}
\hline
\textbf{Инструмент} & \textbf{Назначение в проекте} \\
\hline
TypeScript & Типизация кода на клиенте и сервере для повышения надёжности. \\
\hline
React 19 & Библиотека для построения пользовательского интерфейса. \\
\hline
Vite & Сборка фронтенда (быстрая разработка и оптимизированный билд). \\
\hline
Express & Бэкенд-фреймворк для создания REST API. \\
\hline
SQLite & Лёгкая встраиваемая база данных (используется через better-sqlite3). \\
\hline
Zustand & Управление состоянием на клиенте (глобальные сторы). \\
\hline
TanStack Query & Работа с серверным состоянием (кэширование, синхронизация данных). \\
\hline
Docker & Контейнеризация приложения для простоты развёртывания. \\
\hline
React Hook Form & Управление формами и валидация на клиенте. \\
\hline
Chakra UI & Библиотека компонентов для стилизации интерфейса. \\
\hline
\end{tabularx}
\end{table}

\end{document}